% Discussion, limitations, and proportionate conclusions.
\section{Discussion}\label{sec:discussion}

The generalized gamma kernel-machine promotion-time cure model extends the promotion-time
framework by placing the latent-cause log-intensity in a reproducing-kernel Hilbert space
while retaining a flexible generalized gamma latency and both Poisson and negative-binomial
latent counts. The primary and secondary simulations and the melanoma and breast-cancer
applications evaluate this construction against simpler predictors under a shared protocol.
Their interpretation is limited to the evaluated designs. Kernel families were compared
within prespecified grids, although robustness to broader kernel and hyperparameter choices
remains for future work.

The principal simulation pattern was consistent across the two latent-count laws. On the
quadratic surface, the RBF GG-KM lowered the paired IPCW integrated Brier score against the
linear-kernel GG-KM by about 0.007 to 0.012 and the paired cure root mean squared error by
about 0.04 to 0.07, while its differences from the spline were much smaller and
scenario-dependent. On the periodic surface, the RBF GG-KM improved substantially on the
linear-kernel GG-KM yet trailed the spline by a small margin, and on the linear surface it
paid a small resolved penalty against the correctly specified linear-kernel GG-KM. Within
the original three-surface design, the count law was a secondary factor, since Poisson and
negative-binomial variants within a predictor family differed by at most 0.02 in integrated
Brier score and 0.01 in cure root mean squared error. The melanoma cohort produced overlapping
score distributions in which the linear-kernel Poisson model had the lowest mean integrated
Brier score and GG-KM was competitive with the smallest split-to-split variation. The
breast-cancer cohort showed a complementary pattern: the spline models achieved the lowest
mean integrated Brier score, 0.0877, whereas the mixture-cure model attained the
highest mean cure AUC, 0.7851. Among the GG-KM specifications, the negative-binomial
linear kernel had the lowest mean integrated Brier score, 0.0897, while the exponential
kernel attained the highest mean cure AUC, 0.7124, under both latent-count laws. Together,
the two empirical applications reinforce the simulation message that no
single predictor is uniformly best and that the preferred predictor and
kernel depend on the predictive target and data structure.

Several features of the cure-model setting limit how far these results can be interpreted. Cure fractions are
only weakly identified from right-censored data with finite follow-up, because a long plateau of censored
survivors is consistent with a range of cure probabilities. The melanoma follow-up ends in 1977 and the
estimated cure fractions between 0.577 and 0.648 are model-based extrapolations, not quantities read directly
from the data. The comparison of split-averaged predicted survival with the Kaplan-Meier curve is a
calibration check on the population survival function and does not identify which individuals are cured. The
analysis therefore makes no claim to identify individual cure status, and simulation calibration slopes near
identity do not change that limitation.

Kernel choice and its hyperparameters shape the estimated intensity because they define which subjects are similar
and how rapidly the fitted surface may change. In the primary three-surface comparisons, multiple GG-KM kernel
families were evaluated, with their hyperparameters selected from prespecified grids by leakage-free inner tuning;
the high-dimensional analysis likewise compared multiple kernel families. Bandwidths that are too small can track
idiosyncratic training variation, whereas overly large bandwidths approach a much smoother response; both effects
depend on the training-fold scaling used to measure distance. The high-dimensional results further illustrate that
kernel family can materially affect performance, with polynomial GG-KM performing particularly well on the
interaction surface. A different kernel family or hyperparameter range could therefore shift the balance between
GG-KM and the spline, particularly on surfaces whose geometry does not align well with a given kernel, as the
periodic surface illustrates. Kernel and penalty selection are model-selection steps, and uncertainty from that
selection is not represented by the fitted curve alone. The comparisons fixed a common tuning budget and evaluation
protocol for all methods, so the reported differences describe performance at a shared and modest search intensity
rather than each method's best attainable configuration.

The kernel-machine estimator also carries a computational cost that grows with sample size. An
$n_{\mathrm{train}}\times n_{\mathrm{train}}$ Gram matrix requires quadratic storage, and dense
factorization can require cubic time; nested tuning repeats those operations across candidate kernels and
inner folds. Reusing a Gram matrix within a candidate and retained eigenspace avoids unnecessary work, but
matrices cannot be shared across folds whose preprocessing and training subjects differ. In the primary
simulation experiment, these costs remained manageable at the sample sizes studied, with mean per-fit
runtimes across methods ranging from 1.13 to 40.32 seconds. The richer breast-cancer application was
substantially more computationally demanding, particularly for several nonlinear GG-KM kernels and the
spline models, underscoring that computational burden can increase sharply with a larger and more complex
predictor representation; these timings do not establish scalability to substantially larger cohorts.
Low-rank, Nystr\"om, random-feature, or inducing-point approximations could reduce the burden, although each
changes the effective penalty and out-of-sample predictor and would require separate numerical validation.
Sample size affected the strength of the conclusions as well as runtime: increasing the sample from 300 to
600 reduced Monte Carlo uncertainty and resolved the quadratic and periodic contrasts more clearly.

Censoring and the event definition further bound the interpretation. The evaluation used
inverse-probability-of-censoring weighting with a horizon confined to the region of reliable censoring
support, which controls bias from dependent censoring only to the extent that the censoring model is
adequate. In the melanoma cohort, deaths from other causes were treated as censored. In the breast-cancer
cohort, death from breast cancer was treated as the event, while observations without a recorded
breast-cancer death were censored. Accordingly, both empirical analyses are cause specific in their event
definition and are not competing-risks analyses. Within these latent cause-specific processes,
$p_c(\bm{x})$ describes the probability of no disease-specific death from the modelled process; it is
neither an all-cause survival probability nor a competing-risks cumulative incidence. The interpretation
relies on the corresponding independent-censoring assumption. The negative-binomial dispersion was only
weakly identified from cure and latency information: under Poisson generation the estimated dispersion
correctly stayed near zero, but under true dispersion of 1.00 it was systematically underestimated, ranging
from 0.053 to 0.239. This weak identifiability explains the frequent contact of the dispersion estimate with
its lower bound and cautions against interpreting the fitted dispersion as a precise measure of
overdispersion.

Both empirical applications are observational, and no causal interpretation is attached to the covariate effects
or to the predicted cure fractions. The reported associations describe out-of-sample predictive behaviour under
a fixed analysis protocol, not the effect of intervening on any covariate.

Future work could extend the estimator in directions that were deliberately left outside the present scope.
A controlled comparison against reproducible support-vector or deep cure models would sharpen the
positioning of kernel machines among nonlinear approaches, but a fair comparison requires public
implementations of the same promotion-time model with common generalized gamma latency, which were not
available, so such comparators are noted as open rather than completed. Covariate-dependent latency,
low-rank Gram approximations for larger cohorts, adaptive or multiple-kernel methods, broader kernel
hyperparameter searches, and joint modelling of the censoring mechanism are natural extensions. The current
contribution is the validated combination of an RKHS latent-cause intensity, generalized gamma latency,
Poisson and negative-binomial latent counts, penalized EM and ECME estimation, and reproducible
leakage-free evaluation with paired Monte Carlo uncertainty.
